\chapter{ฟิสิกส์นิวเคลียร์และปฏิกิริยานิวเคลียร์}
\label{chap:nuclear_physics}
\index{ฟิสิกส์นิวเคลียร์}
\index{พลังงานยึดเหนี่ยว}
\index{ฟิชชัน}
\index{ฟิวชัน}

\begin{chapterobjectives}
เมื่อศึกษาบทที่ 5 จบแล้ว นิสิตและผู้ศึกษาสามารถ:
\begin{enumerate}
    \item อธิบายองค์ประกอบของนิวเคลียส คำนวณมวลบกพร่อง (Mass Defect) และพลังงานยึดเหนี่ยวนิวเคลียส (Binding Energy)
    \item วิเคราะห์กราฟพลังงานยึดเหนี่ยวต่อนิวคลีออน ($B/A$) และอธิบายเสถียรภาพของนิวเคลียส
    \item คำนวณสูตรมวลกึ่งเชิงประจักษ์ของไวซ์แซคเกอร์ (Weizsäcker Semi-Empirical Mass Formula)
    \item จำแนกกลไกการสลายตัวของกัมมันตรังสีชนิดแอลฟา บีตา และแกมมา
    \item แก้สมการการสลายตัวลูกโซ่กัมมันตรังสี (Bateman Equations) ด้วยการวิเคราะห์เชิงตัวเลขใน Python
    \item อธิบายเงื่อนไขวิกฤตของปฏิกิริยาลูกโซ่ฟิชชันและสภาวะฟิวชันตามเกณฑ์ของลอว์สัน (Lawson Criterion)
\end{enumerate}
\end{chapterobjectives}

\section{สมบัติของนิวเคลียสและพลังงานยึดเหนี่ยว}
\index{นิวคลีออน}
\index{มวลบกพร่อง}

นิวเคลียสอะตอมประกอบด้วยโปรตอน $Z$ ตัว และนิวตรอน $N$ ตัว รวมกันเป็นเลขมวล $A = Z + N$ รัศมีของนิวเคลียสโดยประมาณคือ:
\begin{equation}
R \approx R_0 A^{1/3}, \quad R_0 \approx 1.2\text{ fm} = 1.2 \times 10^{-15}\text{ m}
\end{equation}

เมื่อนิวคลีออนอิสระมารวมตัวกันเป็นนิวเคลียส มวลของนิวเคลียส $M(A, Z)$ จะน้อยกว่าผลรวมมวลของนิวคลีออนเดี่ยวเสมอ ผลต่างนี้เรียกว่า \textbf{มวลบกพร่อง (Mass Defect, $\Delta m$)}:
\begin{equation}
\Delta m = Z m_p + N m_n - M(A, Z)
\end{equation}
พลังงานที่ปลดปล่อยออกมาในกระบวนการรวมตัวนี้เรียกว่า \textbf{พลังงานยึดเหนี่ยว (Binding Energy, $B$)}:
\begin{equation}
B = \Delta m \cdot c^2 = \left[ Z m_p + (A - Z) m_n - M(A, Z) \right] c^2
\end{equation}

\begin{figure}[H]
\centering
\begin{tikzpicture}[scale=0.95]
    % Axes
    \draw[-{Stealth[scale=1.2]}, slateGrey, thick] (0, 0) -- (11, 0) node[right, font=\footnotesize\bfseries] {เลขมวล $A$};
    \draw[-{Stealth[scale=1.2]}, slateGrey, thick] (0, 0) -- (0, 7.5) node[above, font=\footnotesize\bfseries] {$B/A$ (MeV/nucleon)};
    
    % Grid marks
    \foreach \y in {2, 4, 6, 8} {
        \draw[gray!30, dashed] (0, {0.8*\y}) -- (10.5, {0.8*\y});
        \node at (-0.4, {0.8*\y}) [font=\scriptsize] {\y};
    }
    \foreach \x in {50, 100, 150, 200, 250} {
        \draw[gray!30, dashed] ({0.04*\x}, 0) -- ({0.04*\x}, 6.8);
        \node at ({0.04*\x}, -0.3) [font=\scriptsize] {\x};
    }
    
    % Binding energy curve
    \draw[line width=2.0pt, color=cyberBlue, smooth] plot coordinates {
        (0.16, 1.1) (0.32, 2.8) (0.48, 5.3) (0.64, 7.07) (0.96, 7.4) 
        (1.28, 7.9) (2.24, 8.79) (4.0, 8.6) (6.0, 8.3) (8.0, 7.9) (9.52, 7.57)
    };
    
    % Peak at Iron-56
    \filldraw[quantumGold] (2.24, 7.03) circle (3pt) node[above=4pt, font=\footnotesize\bfseries\color{darkNavy}] {$^{56}\text{Fe}$ ($8.8\text{ MeV}$)};
    
    % Fusion and Fission Regions
    \draw[-{Stealth[scale=1.2]}, emeraldGlow!80!black, line width=1.5pt] (0.8, 4.0) -- (2.0, 5.5);
    \node at (1.2, 3.5) [font=\footnotesize\bfseries\color{emeraldGlow!80!black}] {ปฏิกิริยาฟิวชัน (Fusion)};
    
    \draw[-{Stealth[scale=1.2]}, red!80!black, line width=1.5pt] (9.0, 6.0) -- (7.5, 6.5);
    \node at (8.5, 5.5) [font=\footnotesize\bfseries\color{red!80!black}] {ปฏิกิริยาฟิชชัน (Fission)};
\end{tikzpicture}
\caption{กราฟพลังงานยึดเหนี่ยวนิวเคลียสต่อนิวคลีออน ($B/A$) แสดงความเสถียรสูงสุดที่เหล็ก-56}
\label{fig:binding_energy_curve}
\end{figure}

\section{แบบจำลองหยดของเหลวและสูตรมวลไวซ์แซคเกอร์ (SEMF)}
\index{SEMF}
\index{แบบจำลองหยดของเหลว}

คาร์ล ฟรีดริช ฟอน ไวซ์แซคเกอร์ (Carl Friedrich von Weizsäcker) เสนอสูตรมวลกึ่งเชิงประจักษ์ (Semi-Empirical Mass Formula: SEMF) โดยมองนิวเคลียสเสมือนหยดของเหลวที่มีแรงตึงผิวและประจุไฟฟ้า:
\begin{equation}
B(A, Z) = a_v A - a_s A^{2/3} - a_c \frac{Z(Z-1)}{A^{1/3}} - a_a \frac{(A - 2Z)^2}{A} + \delta(A, Z)
\end{equation}
โดยที่ค่าคงที่จากการทดลองโดยประมาณคือ:
\begin{itemize}
    \item $a_v \approx 15.8\text{ MeV}$ (พจน์ปริมาตร: Volume Term)
    \item $a_s \approx 18.3\text{ MeV}$ (พจน์พื้นที่ผิว: Surface Term)
    \item $a_c \approx 0.714\text{ MeV}$ (พจน์แรงผลักคูลอมบ์: Coulomb Term)
    \item $a_a \approx 23.2\text{ MeV}$ (พจน์อสมมาตร: Asymmetry Term)
    \item $\delta(A, Z)$ คือพจน์การจับคู่ (Pairing Term): เป็น $+a_p A^{-1/2}$ สำหรับคู่-คู่, $0$ สำหรับคี่-คู่, และ $-a_p A^{-1/2}$ สำหรับคี่-คี่
\end{itemize}

\section{กัมมันตภาพรังสีและสมการลูกโซ่ของเบตแมน (Bateman Equations)}
\index{กัมมันตภาพรังสี}
\index{สมการเบตแมน}
\index{ครึ่งชีวิต}

กฎการสลายตัวของสารกัมมันตรังสีชนิดเดียวคือ:
\begin{equation}
N(t) = N_0 e^{-\lambda t} = N_0 2^{-t/t_{1/2}}
\end{equation}
โดยที่ $\lambda = \frac{\ln 2}{t_{1/2}}$ คือค่าคงตัวการสลายตัว (Decay Constant)

ในกรณีการสลายตัวต่อเนื่องเป็นลูกโซ่ ($N_1 \xrightarrow{\lambda_1} N_2 \xrightarrow{\lambda_2} N_3 \xrightarrow{\lambda_3} \dots$):
\begin{equation}
\frac{dN_1}{dt} = -\lambda_1 N_1, \quad \frac{dN_2}{dt} = \lambda_1 N_1 - \lambda_2 N_2, \quad \frac{dN_3}{dt} = \lambda_2 N_2 - \lambda_3 N_3
\end{equation}
แฮร์รี เบตแมน (Harry Bateman) ได้แก้สมการเชิงอนุพันธ์นี้ในรูปผลเฉลยแม่นตรง:
\begin{equation}
N_n(t) = \sum_{i=1}^n c_i e^{-\lambda_i t}, \quad c_i = N_1(0) \left[ \prod_{j=1}^{n-1} \lambda_j \right] \prod_{j=1, j \neq i}^n \frac{1}{\lambda_j - \lambda_i}
\end{equation}

\begin{figure}[H]
\centering
\begin{tikzpicture}[scale=0.95]
    % Axes
    \draw[-{Stealth[scale=1.2]}, slateGrey, thick] (0, 0) -- (8.5, 0) node[right, font=\footnotesize\bfseries] {จำนวนโปรตอน $Z$};
    \draw[-{Stealth[scale=1.2]}, slateGrey, thick] (0, 0) -- (0, 7.5) node[above, font=\footnotesize\bfseries] {จำนวนนิวตรอน $N$};
    
    % N = Z line
    \draw[gray!50, dashed, thick] (0, 0) -- (7.0, 7.0) node[above, font=\scriptsize] {$N = Z$};
    
    % Valley of Stability curve
    \draw[line width=2.2pt, color=darkNavy, smooth] plot coordinates {
        (0, 0) (1.5, 1.5) (3.0, 3.4) (4.5, 5.5) (6.0, 7.2)
    };
    \node at (5.2, 6.2) [font=\footnotesize\bfseries\color{darkNavy}, rotate=48] {แถบเสถียรภาพ (Valley of Stability)};
    
    % Decay modes arrows
    % Beta-minus (N -> P): delta Z = +1, delta N = -1
    \draw[-{Stealth[scale=1.3]}, line width=1.5pt, color=cyberBlue] (3.0, 5.0) -- (3.8, 4.2);
    \node at (2.4, 5.2) [font=\scriptsize\bfseries\color{cyberBlue}, align=center] {สลายตัว $\beta^-$\\($N \to P$)};
    
    % Beta-plus (P -> N): delta Z = -1, delta N = +1
    \draw[-{Stealth[scale=1.3]}, line width=1.5pt, color=neonPurple] (5.2, 3.8) -- (4.4, 4.6);
    \node at (5.9, 3.6) [font=\scriptsize\bfseries\color{neonPurple}, align=center] {สลายตัว $\beta^+$\\($P \to N$)};
    
    % Alpha decay: delta Z = -2, delta N = -2
    \draw[-{Stealth[scale=1.3]}, line width=1.8pt, color=red!80!black] (6.2, 6.8) -- (5.2, 5.8);
    \node at (6.8, 6.0) [font=\scriptsize\bfseries\color{red!80!black}, align=center] {สลายตัว $\alpha$\\($-2P, -2N$)};
    
    % Fission region at high Z, N
    \draw[fill=quantumGold!30, draw=quantumGold, thick, dashed] (5.8, 6.5) circle (0.8cm);
    \node at (5.8, 7.6) [font=\scriptsize\bfseries\color{quantumGold!90!black}] {ย่านเกิดฟิชชัน};
\end{tikzpicture}
\caption{แผนผังนิวไคลด์ของเซเกร (Segrè Chart) แสดงแถบเสถียรภาพและทิศทางการสลายตัวกัมมันตรังสี}
\label{fig:segre_chart}
\end{figure}

\section{ปฏิกิริยานิวเคลียร์ฟิชชันและฟิวชัน}
\index{ฟิชชัน}
\index{ฟิวชัน}
\index{เกณฑ์ของลอว์สัน}

\subsection{นิวเคลียร์ฟิชชัน (Nuclear Fission)}
เมื่อยูเรเนียม-235 จับนิวตรอนช้า (Thermal Neutron) จะเกิดการแตกตัวให้นิวเคลียสขนาดกลางและนิวตรอนความเร็วสูง 2--3 ตัว เช่น:
\begin{equation}
^{235}_{\ 92}\text{U} + ^1_0\text{n} \to ^{236}_{\ 92}\text{U}^* \to ^{141}_{\ 56}\text{Ba} + ^{92}_{36}\text{Kr} + 3 ^1_0\text{n} + Q
\end{equation}
พลังงานปลดปล่อย $Q \approx 200\text{ MeV}$ ต่อการฟิชชันหนึ่งครั้ง

\begin{figure}[H]
\centering
\begin{tikzpicture}[scale=0.92]
    % Step 1: Neutron hit U-235
    \draw[-{Stealth[scale=1.1]}, line width=1.2pt, color=darkNavy] (-4.5, 0) -- (-3.6, 0);
    \filldraw[darkNavy] (-4.5, 0) circle (3pt) node[below=3pt, font=\tiny] {$^1\text{n}$};
    \draw[fill=quantumGold!30, draw=quantumGold!90!black, thick] (-2.8, 0) circle (0.7cm);
    \node at (-2.8, 0) [font=\tiny\bfseries] {$^{235}\text{U}$};
    \node at (-2.8, -1.2) [font=\tiny\bfseries] {1. ดูดกลืนนิวตรอน};
    
    \draw[-{Stealth[scale=1.2]}, slateGrey, thick] (-1.8, 0) -- (-1.2, 0);
    
    % Step 2: Oscillating Compound Nucleus
    \draw[fill=red!20, draw=red!80!black, thick] (0, 0) ellipse (0.95cm and 0.6cm);
    \node at (0, 0) [font=\tiny\bfseries] {$^{236}\text{U}^*$};
    \node at (0, -1.2) [font=\tiny\bfseries] {2. สั่นไหวเสียรูป};
    
    \draw[-{Stealth[scale=1.2]}, slateGrey, thick] (1.2, 0) -- (1.8, 0);
    
    % Step 3: Necking / Dumbbell
    \draw[fill=red!25, draw=red!80!black, thick] (2.7, 0) circle (0.5cm);
    \draw[fill=red!25, draw=red!80!black, thick] (3.9, 0) circle (0.5cm);
    \draw[line width=1.5pt, red!80!black] (3.1, 0.3) -- (3.5, 0.1);
    \draw[line width=1.5pt, red!80!black] (3.1, -0.3) -- (3.5, -0.1);
    \node at (3.3, -1.2) [font=\tiny\bfseries] {3. คอดตัวเป็นดัมเบลล์};
    
    \draw[-{Stealth[scale=1.2]}, slateGrey, thick] (4.6, 0) -- (5.2, 0);
    
    % Step 4: Fission Fragments + Prompt Neutrons
    \draw[fill=cyan!30, draw=cyberBlue, thick] (6.5, 1.0) circle (0.55cm);
    \node at (6.5, 1.0) [font=\tiny\bfseries] {$^{141}\text{Ba}$};
    \draw[fill=emeraldGlow!30, draw=emeraldGlow, thick] (6.5, -1.0) circle (0.45cm);
    \node at (6.5, -1.0) [font=\tiny\bfseries] {$^{92}\text{Kr}$};
    
    % 3 prompt neutrons
    \draw[-{Stealth[scale=1.1]}, line width=1.2pt, color=darkNavy] (6.2, 0.2) -- (7.8, 0.4);
    \filldraw[darkNavy] (7.8, 0.4) circle (2.5pt) node[right, font=\tiny] {$^1\text{n}$};
    \draw[-{Stealth[scale=1.1]}, line width=1.2pt, color=darkNavy] (6.2, 0.0) -- (7.8, 0.0);
    \filldraw[darkNavy] (7.8, 0.0) circle (2.5pt) node[right, font=\tiny] {$^1\text{n}$};
    \draw[-{Stealth[scale=1.1]}, line width=1.2pt, color=darkNavy] (6.2, -0.2) -- (7.8, -0.4);
    \filldraw[darkNavy] (7.8, -0.4) circle (2.5pt) node[right, font=\tiny] {$^1\text{n}$};
    
    \node at (7.0, -1.2) [font=\tiny\bfseries] {4. แตกตัว + $200\text{ MeV}$};
\end{tikzpicture}
\caption{ขั้นตอนการแตกตัวของนิวเคลียสในปฏิกิริยานิวเคลียร์ฟิชชันตามแบบจำลองหยดของเหลว (Liquid Drop Model)}
\label{fig:fission_stages}
\end{figure}

\subsection{นิวเคลียร์ฟิวชัน (Nuclear Fusion)}
ปฏิกิริยารวมตัวของนิวเคลียสเบา เช่น วงจรดิวเทอเรียม-ทริเทียม (D-T):
\begin{equation}
^2_1\text{H} + ^3_1\text{H} \to ^4_2\text{He} + ^1_0\text{n} + 17.6\text{ MeV}
\end{equation}
\textbf{เกณฑ์ของลอว์สัน (Lawson Criterion)} ระบุเงื่อนไขการจุดระเบิดฟิวชันแบบพึ่งพาตัวเอง (Ignition):
\begin{equation}
n \cdot T \cdot \tau_E \ge 3 \times 10^{21}\text{ keV}\cdot\text{s}\cdot\text{m}^{-3}
\end{equation}
โดยที่ $n$ คือความหนาแน่นของพลาสมา, $T$ คืออุณหภูมิ, และ $\tau_E$ คือเวลาการกักเก็บพลังงาน

\section{สถาปัตยกรรมการจำลองเสมือนจริง AR/XR: Nuclear Chain Reaction Simulator}
\index{WebXR}
\index{Fission AR}

\begin{arbox}[การจำลองปฏิกิริยาลูกโซ่นิวเคลียร์ฟิชชัน 3 มิติ แบบ Real-Time]
\textbf{วัตถุประสงค์การเรียนรู้ผ่าน AR:}
\begin{enumerate}
    \item แสดงก้อนมวลยูเรเนียม $^{235}\text{U}$ เสริมสมรรถนะในรูปทรงลูกบาศก์ 3 มิติ ลอยตัวในห้องจริง
    \item เมื่อผู้เรียนยิงนิวตรอนเข้าสู่ก้อนมวล ระบบจะทำการจำลองการแตกตัวแบบ Monte Carlo Particle Cascade แบบ 60 FPS
    \item ผู้เรียนสามารถเลื่อนแท่งควบคุมนิวตรอน (Control Rods) ที่ทำจากแคดเมียมหรือโบรอนเข้า-ออก เพื่อสังเกตค่าตัวคูณการคูณจำนวนนิวตรอนประสิทธิผล $k_{\text{eff}}$:
    \begin{itemize}
        \item $k_{\text{eff}} < 1$ : สภาวะต่ำกว่าวิกฤต (Subcritical) ปฏิกิริยามอดดับลง
        \item $k_{\text{eff}} = 1$ : สภาวะวิกฤต (Critical) เครื่องปฏิกรณ์เดินเครื่องด้วยกำลังคงที่
        \item $k_{\text{eff}} > 1$ : สภาวะเหนือวิกฤต (Supercritical) ปฏิกิริยาทวีคูณอย่างรวดเร็ว
    \end{itemize}
\end{enumerate}
\end{arbox}

\section{การคำนวณเชิงตัวเลขด้วย Python: Bateman Decay Chain Solver}
\index{Python}
\index{Bateman Solver}

\begin{pythonbox}[nuclear\_bateman\_decay.py: การแก้สมการลูกโซ่การสลายตัวกัมมันตรังสี]
import numpy as np
from scipy.integrate import solve_ivp

def bateman_odes(t, N, lambdas):
    """
    ระบบสมการเชิงอนุพันธ์ลูกโซ่กัมมันตรังสี: N1 -> N2 -> N3 (เสถียร)
    """
    l1, l2 = lambdas
    dN1_dt = -l1 * N[0]
    dN2_dt = l1 * N[0] - l2 * N[1]
    dN3_dt = l2 * N[1]
    return [dN1_dt, dN2_dt, dN3_dt]

# กำหนดค่าครึ่งชีวิต
t_half_1 = 5.0   # วัน (ธาตุแม่ A)
t_half_2 = 12.0  # วัน (ธาตุลูก B)
l1 = np.log(2.0) / t_half_1
l2 = np.log(2.0) / t_half_2

N0 = [1.0e6, 0.0, 0.0]  # จำนวนอะตอมเริ่มต้น
t_span = (0.0, 50.0)    # จำลองเป็นเวลา 50 วัน
t_eval = np.linspace(0, 50, 200)

sol = solve_ivp(bateman_odes, t_span, N0, args=([l1, l2],), t_eval=t_eval)

# หาเวลาที่ธาตุลูก B มีปริมาณสูงสุด (Secular / Transient Equilibrium Peak)
max_idx = np.argmax(sol.y[1])
t_peak = sol.t[max_idx]
N2_max = sol.y[1][max_idx]

print("=== Nuclear Decay Chain Simulation ===")
print(f"Decay constant lambda_1: {l1:.4f} day^-1 | lambda_2: {l2:.4f} day^-1")
print(f"Peak daughter (B) activity occurs at t = {t_peak:.2f} days with N_B = {N2_max:.0f} atoms")
\end{pythonbox}

\section{ตัวอย่างการแก้โจทย์เชิงวิเคราะห์ตามระเบียบวิธี P-S-C-T}

\begin{psctexample}[พลังงานปลดปล่อยและการสูญเสียมวลในปฏิกิริยาฟิชชันของ U-235]
\textbf{โจทย์ (Problem):} \\
พิจารณาปฏิกิริยาฟิชชันของ $^{235}\text{U}$ โดยนิวตรอนพลังงานต่ำ:
\begin{equation}
^{235}_{\ 92}\text{U} + ^1_0\text{n} \to ^{141}_{\ 56}\text{Ba} + ^{92}_{36}\text{Kr} + 3 ^1_0\text{n}
\end{equation}
กำหนดมวลอะตอมในหน่วยมวลอะตอม ($u$):
\begin{itemize}
    \item $M(^{235}\text{U}) = 235.0439299\text{ u}$
    \item $M(^1\text{n}) = 1.0086649\text{ u}$
    \item $M(^{141}\text{Ba}) = 140.914411\text{ u}$
    \item $M(^{92}\text{Kr}) = 91.926156\text{ u}$
\end{itemize}
(กำหนดค่าสมมูลมวล-พลังงาน: $1\text{ u} \cdot c^2 = 931.494\text{ MeV}$)
\begin{enumerate}
    \item จงคำนวณมวลบกพร่องของปฏิกิริยานี้ในหน่วย $u$
    \item จงหาพลังงานปลดปล่อย $Q$ ในหน่วย $\text{MeV}$
    \item หากโรงไฟฟ้านิวเคลียร์ใช้ $^{235}\text{U}$ บริสุทธิ์ $1.0\text{ kg}$ เกิดฟิชชันสมบูรณ์ จะปลดปล่อยพลังงานทั้งหมดกี่เมกะวัตต์-วัน (MW$\cdot$day)?
\end{enumerate}

\textbf{ยุทธวิธี (Strategy):} \\
\begin{itemize}
    \item มวลบกพร่อง $\Delta M = M_{\text{ตั้งต้น}} - M_{\text{ผลิตภัณฑ์}}$
    \item พลังงานปลดปล่อย $Q = \Delta M \times 931.494\text{ MeV}$
    \item จำนวนนิวเคลียสใน $1.0\text{ kg}$ คำนวณจาก $N = \frac{m}{M_{\text{mol}}} N_A$
\end{itemize}

\textbf{การคำนวณ (Computation):} \\
\textit{1. คำนวณมวลตั้งต้นและมวลผลผลิต:}
\begin{equation}
M_{\text{in}} = M(^{235}\text{U}) + M(^1\text{n}) = 235.043930 + 1.008665 = 236.052595\text{ u}
\end{equation}
\begin{align}
M_{\text{out}} &= M(^{141}\text{Ba}) + M(^{92}\text{Kr}) + 3 M(^1\text{n}) \nonumber \\
&= 140.914411 + 91.926156 + 3(1.008665) = 235.866562\text{ u}
\end{align}
ผลต่างมวลคือ:
\begin{equation}
\Delta M = 236.052595 - 235.866562 \approx 0.186033\text{ u}
\end{equation}

\textit{2. คำนวณค่า $Q$:}
\begin{equation}
Q = 0.186033\text{ u} \times 931.494\text{ MeV/u} \approx 173.29\text{ MeV}
\end{equation}
(หากรวมพลังงานจากการสลายตัวบีตาของผลผลิตฟิชชันเพิ่มเติม จะได้ประมาณ $200\text{ MeV}$)

\textit{3. พลังงานจาก $^{235}\text{U}$ ปริมาณ $1.0\text{ kg}$:}
\begin{equation}
N = \frac{1000\text{ g}}{235\text{ g/mol}} \times 6.022 \times 10^{23}\text{ nuclei/mol} \approx 2.563 \times 10^{24}\text{ nuclei}
\end{equation}
พลังงานรวมในหน่วยจูล:
\begin{equation}
E_{\text{total}} = N \times (173.29\text{ MeV}) \times (1.602 \times 10^{-13}\text{ J/MeV}) \approx 7.115 \times 10^{13}\text{ J}
\end{equation}
แปลงเป็นหน่วยเมกะวัตต์-วัน ($1\text{ MW}\cdot\text{day} = 10^6\text{ W} \times 86400\text{ s} = 8.64 \times 10^{10}\text{ J}$):
\begin{equation}
E = \frac{7.115 \times 10^{13}\text{ J}}{8.64 \times 10^{10}\text{ J/(MW}\cdot\text{day)}} \approx 823.5\text{ MW}\cdot\text{day}
\end{equation}

\textbf{ข้อคิดทางฟิสิกส์ (Takeaway):} \\
ยูเรเนียมเพียง $1\text{ kg}$ ปลดปล่อยพลังงานความร้อนเทียบเท่ากับการเผาไหม้ถ่านหินคุณภาพสูงมากกว่า $2500\text{ ตัน}$ นี่คือพลังอันมหาศาลของแรงนิวเคลียร์อย่างเข้มเมื่อเปรียบเทียบกับปฏิกิริยาเคมีดั้งเดิม
\qedexam
\end{psctexample}

\begin{psctexample}[การคำนวณค่า $Q$-Value และการกระจายพลังงานในปฏิกิริยาฟิวชัน D-T (D-T Fusion)]
ปฏิกิริยาฟิวชันเทอร์โมนิวเคลียร์หลักที่ใช้ในเครื่องปฏิกรณ์นิวเคลียร์ฟิวชันระดับนานาชาติ ITER คือปฏิกิริยาระหว่างดิวเทอเรียม ($^2_1\text{H}$) และทริเทียม ($^3_1\text{H}$):
\begin{equation}
^2_1\text{H} + ^3_1\text{H} \to ^4_2\text{He} + ^1_0\text{n}
\end{equation}
กำหนดมวลนิวเคลียส: $m(^2\text{H}) = 2.014102\text{ u}$, $m(^3\text{H}) = 3.016049\text{ u}$, $m(^4\text{He}) = 4.002603\text{ u}$, $m_n = 1.008665\text{ u}$ และ $1\text{ u} \cdot c^2 = 931.494\text{ MeV}$
\begin{enumerate}
    \item จงคำนวณค่าพลังงานปลดปล่อยสุทธิ $Q$-Value ของปฏิกิริยาในหน่วย $\text{MeV}$
    \item ในกรอบที่อนุภาคตั้งต้นเคลื่อนที่ช้ามากจนโมเมนตัมรวมของระบบเป็นศูนย์ จงคำนวณการแบ่งสันปันส่วนพลังงานจลน์ให้กับอนุภาคแอลฟา ($K_\alpha$) และนิวตรอน ($K_n$)
\end{enumerate}

\textbf{โจทย์ (Problem):} \\
คำนวณพลังงานรวมที่ปลดปล่อยและสัดส่วนพลังงานจลน์ของแต่ละอนุภาคหลังเกิดปฏิกิริยาฟิวชัน

\textbf{ยุทธวิธี (Strategy):} \\
\begin{itemize}
    \item มวลพร่องเชิงปฏิกิริยา: $\Delta m = [m(^2\text{H}) + m(^3\text{H})] - [m(^4\text{He}) + m_n]$
    \item ค่า $Q = \Delta m \cdot c^2$
    \item กฎอนุรักษ์โมเมนตัม: $\mathbf{p}_\alpha + \mathbf{p}_n = 0 \implies p_\alpha = p_n = p$
    \item พลังงานจลน์แบบไม่สัมพัทธภาพ: $K = \frac{p^2}{2m}$ ดังนั้น $\frac{K_\alpha}{K_n} = \frac{m_n}{m_\alpha} \approx \frac{1}{4}$
    \item $K_\alpha + K_n = Q \implies K_n = \frac{m_\alpha}{m_\alpha + m_n} Q, \quad K_\alpha = \frac{m_n}{m_\alpha + m_n} Q$
\end{itemize}

\textbf{การคำนวณ (Computation):} \\
\textit{1. คำนวณมวลพร่องและ $Q$-Value:}
\begin{align}
m_{\text{reactants}} &= 2.014102\text{ u} + 3.016049\text{ u} = 5.030151\text{ u} \nonumber \\
m_{\text{products}} &= 4.002603\text{ u} + 1.008665\text{ u} = 5.011268\text{ u} \nonumber \\
\Delta m &= 5.030151\text{ u} - 5.011268\text{ u} = 0.018883\text{ u}
\end{align}
คำนวณค่า $Q$:
\begin{equation}
Q = (0.018883\text{ u}) \times (931.494\text{ MeV/u}) \approx 17.589\text{ MeV} \approx 17.59\text{ MeV}
\end{equation}

\textit{2. การแบ่งพลังงานจลน์:}
\begin{align}
K_n &= \left(\frac{4.002603}{4.002603 + 1.008665}\right) \times 17.589\text{ MeV} \nonumber \\
&= \left(\frac{4.002603}{5.011268}\right) \times 17.589\text{ MeV} \approx 0.7987 \times 17.589\text{ MeV} \approx 14.05\text{ MeV}
\end{align}
และสำหรับอนุภาคแอลฟา:
\begin{equation}
K_\alpha = Q - K_n = 17.589 - 14.048 \approx 3.54\text{ MeV}
\end{equation}

\textbf{ข้อคิดทางฟิสิกส์ (Takeaway):} \\
พลังงานร้อยละ 80 ของปฏิกิริยา ($14.05\text{ MeV}$) ตกไปอยู่นิวตรอนที่ไม่มีประจุไฟฟ้า นิวตรอนเหล่านี้จะทะลุผ่านสนามแม่เหล็กกักกันพลาสมาออกไปชนกับผนังผ้าห่ม (Blanket) เพื่อเปลี่ยนพลังงานจลน์เป็นความร้อนสำหรับผลิตไอน้ำปั่นกังหัน ขณะที่อนุภาคแอลฟาประจุ $+2$ จะถูกกักไว้ในสนามแม่เหล็กเพื่อคอยถ่ายทอดความร้อนกลับเข้าสู่พลาสมา (Alpha Heating) เพื่อรักษาอุณหภูมิฟิวชันให้คงอยู่ต่อไป
\qedexam
\end{psctexample}

\section{สรุปสาระสำคัญประจำบท}

\begin{table}[H]
\centering
\small
\begin{tabularx}{\textwidth}{l >{\centering\arraybackslash}p{4.2cm} Y}
\toprule
\textbf{กระบวนการ} & \textbf{สมการ / สูตร} & \textbf{ความสำคัญทางฟิสิกส์} \\
\midrule
พลังงานยึดเหนี่ยว & $B = \Delta m \cdot c^2$ & มาตรวัดเสถียรภาพของนิวเคลียสอะตอม \\
SEMF & $B \approx a_v A - a_s A^{2/3} - \dots$ & อธิบายสมบัติมวลนิวเคลียสด้วยแบบจำลองหยดของเหลว \\
การสลายตัวกัมมันตรังสี & $N(t) = N_0 e^{-\lambda t}$ & พลวัตเชิงสถิติของการเปลี่ยนรูปนิวเคลียสที่ไม่เสถียร \\
ฟิชชันนิวเคลียร์ & $^{235}\text{U} + \text{n} \to \text{Products} + 200\text{ MeV}$ & การแตกตัวของนิวเคลียสหนักปลดปล่อยพลังงานสูง \\
ฟิวชันนิวเคลียร์ & $^2\text{H} + ^3\text{H} \to ^4\text{He} + \text{n} + 17.6\text{ MeV}$ & แหล่งกำเนิดพลังงานของดาวฤกษ์และเตาปฏิกรณ์ในอนาคต \\
\bottomrule
\end{tabularx}
\caption{สรุปหลักการสำคัญในฟิสิกส์นิวเคลียร์และปฏิกิริยานิวเคลียร์}
\label{tab:nuclear_summary}
\end{table}

\section{แบบฝึกหัดท้ายบท}
\begin{enumerate}
    \item จงคำนวณพลังงานยึดเหนี่ยวนิวเคลียสต่อนิวคลีออนของ $^{4}_{2}\text{He}$ กำหนดมวลอะตอม $M(^4\text{He}) = 4.002603\text{ u}$, $m_p = 1.007276\text{ u}$, $m_n = 1.008665\text{ u}$
    \item ไอโซโทปกัมมันตรังสีชนิดหนึ่งมีกิจกรรมกัมมันตรังสีลดลงเหลือ $12.5\%$ ของค่าเริ่มต้นในเวลา 18 วัน จงหาครึ่งชีวิตและค่าคงตัวการสลายตัวของไอโซโทปนี้
    \item จงอธิบายบทบาทของสารหน่วงนิวตรอน (Moderator) และแท่งควบคุม (Control Rods) ในเครื่องปฏิกรณ์นิวเคลียร์ฟิชชัน
    \item ไอโซโทปคาร์บอน-14 ($^{14}_{6}\text{C}$) สลายตัวให้อนุภาคบีตาลบด้วยครึ่งชีวิต $T_{1/2} = 5730\text{ ปี}$ ตัวอย่างไม้โบราณที่ขุดพบมีกิจกรรมกัมมันตรังสีของ $^{14}\text{C}$ เพียง $25.0\%$ เมื่อเทียบกับไม้มีชีวิตในปัจจุบัน จงประเมินอายุของโบราณวัตถุชิ้นนี้
    \item จงอธิบายเกณฑ์ของลอว์สัน (Lawson Criterion: $n \tau_E \ge 10^{20}\text{ m}^{-3}\text{s}$) สำหรับการจุดระเบิดพลาสมาฟิวชัน (Fusion Ignition) และเปรียบเทียบการกักขังพลาสมาด้วยสนามแม่เหล็ก (Magnetic Confinement: Tokamak) กับการกักขังด้วยความเฉื่อย (Inertial Confinement: Laser Fusion)
\end{enumerate}
